
\documentclass[12pt]{article}
\usepackage[margin=1in]{geometry}
\usepackage{kotex}
\usepackage{setspace} \doublespacing
\title{마지막 남은 비대칭성 II: 공정한 사회일수록 \\왜 젊은 여성은 더 진보적이 되는가}
\author{임강빈\thanks{Email: yim@sch.ac.kr (순천향대학교 교수)}}
\date{2026년 6월}
\begin{document}
\maketitle

최근 세계 여러 나라에서는 흥미로운 현상이 관찰되고 있다. 젊은 여성은 점점 더 진보적인 정치 성향을 보이는 반면, 젊은 남성은 상대적으로 보수적인 성향을 보이는 경우가 늘어나고 있다. 이러한 현상은 특정 국가만의 이야기가 아니다. 북미와 서유럽, 북유럽, 그리고 동아시아 일부 국가들에서도 비슷한 경향이 보고되고 있다. 특히 한국은 젊은 남성과 젊은 여성의 정치적 성향 차이가 가장 크게 나타나는 국가 중 하나로 자주 언급된다. 많은 사람들은 이러한 현상을 정치적 선전, 세대 갈등, 또는 젠더 갈등의 결과로 설명하려고 한다. 그러나 정말 그것이 본질일까? 정치적 성향의 변화는 정치 자체에서 비롯된 것일까? 아니면 그보다 더 깊은 사회 구조의 변화가 정치적 성향의 변화로 나타나고 있는 것은 아닐까? 이 질문에 답하기 위해서는 먼저 현대 사회가 지난 수십 년 동안 어떤 방향으로 변화해 왔는지를 살펴볼 필요가 있다.

인류 문명의 역사는 여러 측면에서 불평등의 감소 과정이었다. 왕과 귀족만이 정치에 참여하던 시대가 있었다.
교육은 일부 계층만의 특권이었다. 여성은 투표권을 갖지 못했고, 재산권도 제한되었다. 출신 계급은 개인의 삶을 결정하는 중요한 요소였다. 그러나 현대 사회는 이러한 장벽들을 지속적으로 제거해 왔다. 정치 참여는 확대되었다. 교육 기회는 보편화되었다. 경제 활동의 문턱은 낮아졌다. 여성의 사회 참여는 크게 증가했다.
오늘날 선진국이라 불리는 국가들은 역사상 가장 평등하고 가장 개방적이며 가장 공정한 사회에 가깝다.

그런데 바로 여기서 흥미로운 역설이 등장한다. 공정성이 확대될수록 사람들은 오히려 남아 있는 불공정을 더욱 민감하게 인식하기 시작한다는 것이다. 더러운 유리창에 묻은 작은 얼룩은 눈에 잘 띄지 않는다. 그러나 깨끗하게 닦인 유리창 위의 작은 얼룩은 누구나 쉽게 발견할 수 있다. 사회도 마찬가지다. 많은 차별이 제거되면 남아 있는 차별은 더욱 선명하게 보인다. 많은 불평등이 해소되면 남아 있는 불평등은 더욱 크게 느껴진다.

실제로 역사를 돌아보면 이러한 현상은 반복적으로 나타난다. 노예제가 일반적이던 시대에는 노예제 자체가 문제로 인식되지 않았다. 그러나 노예제가 약화되기 시작하자 남아 있는 노예제는 더욱 부당하게 보이기 시작했다. 인종차별이 일상이던 시대에는 차별 자체가 자연스럽게 받아들여졌다. 그러나 법적 차별이 사라질수록 남아 있는 차별은 더욱 예민하게 인식되었다. 여성의 사회적 권리가 확대되면서도 비슷한 일이 벌어졌다. 교육과 취업, 정치 참여에서의 장벽이 제거될수록 남아 있는 불평등에 대한 관심은 오히려 더욱 커졌다. 진보는 불평등에 대한 인식을 없애는 것이 아니라, 오히려 더욱 선명하게 만드는 경우가 많다.

이러한 관점에서 최근의 정치적 변화도 다시 볼 수 있다. 현대 사회는 남성과 여성을 동일한 경쟁 구조 안으로 편입시켰다. 남성과 여성은 같은 대학에 진학한다. 같은 시험을 치른다. 같은 직장을 목표로 한다. 같은 조직에서 승진을 경쟁한다. 같은 정치적 권리를 가진다. 즉 사회적 역할의 상당 부분이 대칭화되었다. 이것은 현대 문명이 이룬 가장 중요한 성취 중 하나다. 그런데 바로 그 순간 새로운 질문이 등장한다. 정말 모든 것이 대칭적인가?

많은 영역에서 남녀의 차이는 줄어들었지만, 한 가지는 여전히 남아 있다. 바로 임신과 출산이다. 임신은 여성의 몸에서 이루어진다. 출산도 여성의 몸에서 이루어진다. 출산 이후의 회복 역시 여성의 몸에서 이루어진다. 기술은 부담을 줄일 수 있다. 정책은 비용을 보전할 수 있다. 그러나 생물학적 현실 자체를 제거할 수는 없다. 그리고 사회적 역할이 대칭화될수록 이 차이는 더욱 선명하게 보이기 시작한다. 과거에는 수많은 역할 차이 중 하나였던 것이 이제는 거의 유일하게 남은 비대칭성처럼 보이게 된다.

바로 이 지점이 출산율 문제와 정치적 성향 변화를 동시에 설명할 수 있는 중요한 단서일 수 있다. 만약 한 사회가 공정성을 핵심 가치로 삼고 있다면, 사람들은 남아 있는 비대칭성에 더욱 민감해질 수밖에 없다. 그리고 그 비대칭성이 생애 전체에 영향을 미치는 문제라면 더욱 그렇다. 왜 동일한 경쟁을 요구받는데 생물학적 부담은 동일하지 않은가? 왜 교육과 직업, 승진에서는 평등을 이야기하면서 출산은 여전히 특정 성별의 몸에 의존하는가? 왜 현대 사회가 만들어 낸 공정 경쟁 구조 안에서 출산은 여전히 예외로 남아 있는가?
이 질문들은 단순한 경제 문제가 아니다. 주거비나 양육비만으로 설명되지 않는다. 그것은 공정성에 대한 질문이다.

흥미로운 점은 이러한 현상이 선진국일수록 더욱 강하게 나타난다는 것이다. 국제 비교 연구들은 여성의 진보 성향이 성평등 수준이 높은 국가들에서 더 뚜렷하게 나타나는 경향을 보고하고 있다. 또한 젊은 세대에서 이러한 현상이 더욱 강하게 나타난다는 결과도 반복적으로 발견된다. 물론 이것이 모든 국가에서 동일하게 나타나는 것은 아니다. 그러나 적어도 한 가지는 분명해 보인다. 현대 사회가 공정성을 확대할수록 사람들은 남아 있는 비대칭성을 더욱 민감하게 인식하게 된다는 사실이다.

출산율 저하와 젊은 여성의 진보화는 흔히 별개의 문제로 다루어진다. 하나는 인구 문제이고, 다른 하나는 정치 문제로 취급된다. 그러나 어쩌면 두 현상은 같은 구조적 변화의 서로 다른 표현일 수 있다. 사회적 평등이 확대될수록 남아 있는 차이는 더욱 선명해진다. 그리고 그 차이에 대한 문제의식은 더 많은 공정성을 요구하게 만든다. 따라서 오늘날 우리가 목격하는 정치적 변화는 단순한 이념의 이동이 아닐 수도 있다. 그것은 현대 사회가 자신의 성공 이후에 마주한 새로운 질문일 수 있다. 과거의 불평등이 대부분 제거된 이후에도 여전히 남아 있는 차이를 우리는 어떻게 이해해야 하는가. 그리고 그 차이를 안고 살아가는 사회는 어떤 방향으로 변화하게 될 것인가. 어쩌면 이것이 21세기 선진국들이 공통적으로 마주하고 있는 가장 중요한 질문 중 하나인지도 모른다.

\end{document}
